% !PRINTED 440
\weberpage{440}
%<<CONT
\gesperrt{Die Gesammtheit der Zahlen $(\varrho)$ stimmt also, vom Vorzeichen und von der Reihenfolge abgesehen, mit der Gesammtheit der Zahlen $(\nu)$ überein.}

In den Formeln (1) und (2) können wir, aber wegen der Periodicität von Sinus und Cosinus, $\nu m$ durch $\varrho$ ersetzen, und da $\cos(-\varphi) = \cos\varphi$ ist, so können wir in der Formel (1) $m\nu$ auch durch $\nu$ ersetzen, d. h. das Vorzeichen ist von $m$ nicht abhängig. Um es zu bestimmen, berücksichtigen wir, dass der Cosinus eines Winkels im ersten Quadranten positiv, im zweiten Quadranten negativ ist, dass also das Vorzeichen in (1) positiv oder negativ ist, je nachdem von den Winkeln
\[
\frac{2\nu\pi}{n}
\]
eine gerade oder eine ungerade Anzahl im zweiten Quadranten liegt, oder je nachdem eine gerade oder eine ungerade Anzahl von Zahlen $\nu$ zwischen $\dfrac{n}{4}$ und $\dfrac{n}{2}$ liegt.

Wir bezeichnen, wenn $x$ irgend eine nicht ganze Zahl ist, mit $E(x)$ die grösste ganze Zahl, die kleiner als $x$ ist, so dass $x$ zwischen $E(x)$ und $E(x)+1$ liegt. Die Anzahl der ganzen Zahlen, die zwischen zwei Zahlen $x$ und $y$ liegt, ist dann gleich $E(y) - E(x)$, und wir haben also zu untersuchen, ob
\[
E\left(\frac{n}{2}\right) - E\left(\frac{n}{4}\right)
\]
eine gerade oder eine ungerade Zahl ist. Wir müssen vier Fälle unterscheiden, wie sich aus der folgenden Zusammenstellung ergiebt, worin $k$ eine nicht negative ganze Zahl bedeutet:
\[
\begin{array}{lll}
n = 8k+1, & E\left(\dfrac{n}{2}\right) = 4k, & E\left(\dfrac{n}{4}\right) = 2k, \\[2ex]
n = 8k+3, & E\left(\dfrac{n}{2}\right) = 4k+1, & E\left(\dfrac{n}{4}\right) = 2k, \\[2ex]
n = 8k+5, & E\left(\dfrac{n}{2}\right) = 4k+2, & E\left(\dfrac{n}{4}\right) = 2k+1, \\[2ex]
n = 8k+7, & E\left(\dfrac{n}{2}\right) = 4k+3, & E\left(\dfrac{n}{4}\right) = 2k+1.
\end{array}
\]
Es ist also in (1) das positive Zeichen zu nehmen, wenn $n$ von der Form $8k+1$ oder $8k+7$, das negative, wenn $n$ von der Form $8k+3$ oder $8k+5$ ist. In den ersten Fällen ist
%CONT>>
